\cotefile{157-3}
\scanpage{237}
\mspage{187}

\[
\operatorname{Hom}_{\pi}(\pi_{n-1}(Z,x),\underset{\substack{\Vert\\\inserted{F_{\unsure{x|n}}}}}{\deleted{\unsure{A}}M})=0
\]

et ceci \uline{\unsure{quel|quels} que \unsure{soit|soient} \unsure{le|les}} $\pi$-\unsure{module|modules} $M$, puisque $\pi=\pi_1(X,x)\simeq\pi_1(Y,y)$. Cela implique (en prenant p.~ex. $M=\pi_{n-1}(Z,x)$ !) que $\pi_{n-1}(Z,x)=0$. (\deleted{\unsure{Il}}Mais en faisant attention dans le cas $n=2$, qu'on sait déjà que $\pi_1(Z,\unsure{x})$ est abélien).

Donc le th. rectifié \unsure{page} $178$ a bien l'air vrai. Il est un peu triste, puisque dans chacune des implications

\[
a_n\Longrightarrow b_n\Longrightarrow c_n\Longrightarrow a_{n-1}
\]

on perd un chouïa, et en fin de compte on sort avec les magnifiques implications

\[
a_n\Longrightarrow a_{n-1}.
\]

De toutes façons, il est $\pm$ clair que le ``bon'' localiseur $W_n$, celui qui doit définir les ``$n$-types d'homotopie'' comme objets de $\operatorname{Hot}_{W_n}$, est défini par la condition $\unsure{a_n}$.

2. \uline{Critère de $W_\infty$-asphéricité}

Je vais essayer de faire mieux quand $Y$ est réduit à un point. Cette fois, je
% archive sheet number: 237
\scanpage{238}
\mspage{188}

veux être \unsure{strict}, et donner (si je peux) des démonstrations complètes dans le cadre où je me place à présent, celui des modéliseurs $\mathrm{Cat}$, et en n'utilisant que ce que je sais y démontrer directement à présent, sans me référer à \deleted{\unsure{les}} la théorie $\frac12$ simpliciale ou topologique.

\uline{Théorème} 2 (?) Soit $X$ dans $\mathrm{Cat}$, $n\in\unsure{\mathbb N|\mathcal N}$ un entier $\geqslant 1$. Alors les conditions suivantes sont équivalentes :

$a_n)$ $X$ est $0$-connexe, et \deleted{\unsure{pour}} si $x$ est dans $X$, $\pi_{\unsure{i|n}}(X,x)=0$ pour $i\in[1,n]$ (condition qui ne dépend pas du choix de $x$), i.e. les \inserted{$\Omega^i(X,x)$ pour $i\in[1,n]$ sont $0$-connexes,}

$b_n)$ Pour \deleted{tout} \deleted{\unsure{$F$ faisc. constant}} \deleted{\unsure{faisc.}} groupe \deleted{\unsure{abélien}} $G$, $G\longrightarrow H^0(X,G_X)$ est un iso, et \deleted{\unsure{pour}} $H^1(X,G_X)=0$. De plus si $G$ est abélien, on a aussi $H^i(X,G_X)=0$ pour $i\in[1,n]$.

\uline{Rém.} Pour $n=1$ cette équivalence est $\pm$ tautologique et a été rappelée \inserted{(sous une forme légèrement différente)} au début. (p.~71). On procède par récurrence. On peut donc supposer que \inserted{$n\geqslant 2$, que} $X$ satisfait $a_{n-1},b_{n-1}$, et prouver qu'alors $a_n,b_n$ sont \unsure{équivalentes|équivalents}, i.e. $\pi_n(X,x)=0$ équivaut à $H^n(X,\deleted{\unsure{M}}A)=0$ pour tt gr. ab. $\deleted{\unsure{M}}A$.
% archive sheet number: 238
\scanpage{239}
\mspage{189}

On utilise la \deleted{fibration} $W_\infty$-fibration (parfaite)

\diagram{\begin{tikzcd}[ampersand replacement=\&,column sep=0.25em,row sep=1.2em]
E\arrow[d] \& {}=\underline{\mathrm{Ch}}(X,x,{-}) \& \qquad E\ \text{$W_\infty$-asphérique}\\
X \& \&
\end{tikzcd}}

de fibre \inserted{$Z$} en $x$ égale à $\Omega(X,x)$.

Supposons $a_n)$ pour $X$, donc $a_{n-1})$ pour $\Omega(X,x)$. Par hyp. de récurrence, on a donc

\[
H^{\unsure{i|n}}(Z,\deleted{\unsure{A}}M)=0\qquad\text{pour }\deleted{\unsure{q}}\,i\in[1,n-1].
\]

$\deleted{\unsure{A}}M$ \deleted{faisceau constant sur $Z$} groupe abélien. Donc dans la suite spectrale de Leray

\[
H^{\unsure{*|n}}(E,M)\Longleftarrow E^{pq}_2=H^p(X,H^q(Z,M))
\]

on a

\[
E^{pq}_2=0\quad\text{si}\quad q\in[1,n-1]
\]

\diagram{\begin{tikzcd}[ampersand replacement=\&,column sep=.36em,row sep=.24em]
q \& \& \& \& \&\\
|[label=left:{$\scriptstyle n$}]| \bullet\arrow[u] \& \bullet \& \bullet \& \bullet \& \bullet \& \cdots\\
{} \& 0 \& \& 0 \& \&\\
\bullet\arrow[uu,no head]\arrow[rrrrr] \& \bullet \& \bullet \& \bullet \& \bullet \& {} \\
\& \& \& \& \& p

\end{tikzcd}}

On a donc la suite exacte

\[
0\longrightarrow H^n(X,M)\longrightarrow H^n(E,M)
\]

et comme $E$ est $W_\infty$-asphérique, on a $H^n(E,M)=0$, donc $H^n(X,M)=0$.

Inversement, supposons $b_n)$, prouvons $a_n$. On sait déjà qu'on a $a_{n-1}$, donc $a_{n-2}$ pour $\Omega(X,x)$, \deleted{\unsure{i.e.}} qui implique

\[
H^i(Z,M)=0\qquad\text{pour}\qquad i\in[1,n-2],
\]

\[
\text{donc}\qquad E^{pq}_2=0\quad\text{si}\quad q\in[1,n-2].
\]

On a donc la suite exacte
% archive sheet number: 239
\scanpage{240}
\mspage{190}

\diagram{\begingroup\scriptsize
\begin{tikzcd}[ampersand replacement=\&,column sep=.65em,row sep=.35em]
0\arrow[r,no head] \& H^{n-1}(X,M)\arrow[r,no head]\arrow[d,equal] \& H^{n-1}(E,M)\arrow[r,no head]\arrow[d,equal] \& H^0(X,H^{n-1}(Z,M))\arrow[r] \& H^n(X,M)\arrow[r]\arrow[d,equal] \& H^n(E\unsure{,|),}M
\\
\&0\&0\&\&0\&
\end{tikzcd}\endgroup}

\[
\text{donc}\qquad H^0(X,\underbrace{H^{n-1}(Z,M)}_{\substack{\text{système local}\\\text{\uline{constant} sur }X\\\text{car }X\text{ est }1\text{-connexe}}})=0
\]

\[
\text{\unsure{donc}}\qquad H^{n-1}(Z,M)=0
\]

ce qui, par hyp. de récurrence, implique $\pi_{n-1}(Z,\unsure{x})=0$, donc $\pi_n(X,x)=0$, qed.

\uline{Corollaire} Pour que $X$ soit $W_\infty$-asphérique, \unsure{i.e.} pour qu'on ait \deleted{\unsure{pour}} $G\xrightarrow{\sim}H^0(X,G)$, $H^1(\unsure{X},G)=0$ tout groupe, et $H^i(X,G)=0$ si $i\geqslant 2$ pour $G$ abélien, il f. et s. que \deleted{\unsure{les}}

$1^\circ)$ $X$ est $\unsure{0{=}|0\text{-}}$connexe, et ($x$ étant un point de $X$) $2^\circ)$ les $\Omega^i(X,x)$ $0$-connexes pour $i\geqslant 1$.
% archive sheet number: 240
\scanpage{241}
\mspage{191}

3. \uline{Sur la philosophie des fibres homotopiques.}

\deleted{\unsure{3.}} Le critère d'asphéricité (pour $W_\infty$) que je viens d'obtenir, grâce aux foncteurs $\underline{\mathrm{Ch}}(X)$ (ou au choix, $\underline{\mathrm{Ch}}_\infty(X)$) et ses dérivés $\Omega(X,x)$ et les itérés, enlève tout doute (s'il pouvait en rester un) que les \deleted{\unsure{esp.}} ``espaces de chemins'' que j'ai construit (assez péniblement !) sont bel et bien ``les bons''.

Je voudrais maintenant me rendre compte de façon la plus précise possible \deleted{\unsure{des}} de leur ``canonicité'', dans quel mesure ces constructions sont indépendantes d'un choix (p.~ex. celui entre \unsure{la|le} \unsure{théorie|choix} $\underline{\mathrm{Ch}}$, et $\underline{\mathrm{Ch}}_\infty$).

Tout revient, finalement, à des questions de factorisations d'une flèche quelconque

\[
f:X\longrightarrow Y
\]

de $\mathrm{Cat}$ en

\diagram{\begin{tikzcd}[ampersand replacement=\&]
X\arrow[r,"{\deleted{\unsure{f}}i_{\deleted{\unsure{f}} \unsure{f}}}"]\&\overline X\arrow[r,"{\overline f}"]\&Y,
\end{tikzcd}}

avec

\[
\unsure{i_f}\in W,\quad \overline f\in W\text{-}\mathrm{fib}
\]

(le fait que $\overline f$ puisse $=$ être choisie $W_\infty$-parfaite ne me paraissant pas essentiel ici).

\unsure{C'est} bien de cela qu'il s'agit pour construire
% archive sheet number: 241
\scanpage{242}
\mspage{192}

un objet ``Chemins dans $X$''

\[
\underline{\mathrm{Ch}}(X),
\]

en factorisant l'application diagonale

\[
X\xrightarrow{\mathrm{diag}_X}X\times X,
\]

et aussi quand on construit $\Omega(X,x)$ comme fibre en $x$ d'une \inserted{$W_{\unsure{\infty}}$-}fibration

\diagram{\begin{tikzcd}[ampersand replacement=\&]
E(X,x)\arrow[r,very thick]\&(X,x)
\end{tikzcd}}

qui factorise l'inclusion

\[
\{x\}\longrightarrow X
\]

(donc $E(X,x)$ étant $W$-asphérique).

On peut être plus ou moins exigeant et regarder, soit pour un \inserted{``but''} $\unsure{Y}$ (ou plutôt un $S$) fixé dans $\mathrm{Cat}$, une factorisation fonctorielle

\[
C:\mathrm{Cat}_{/S}\longrightarrow\deleted{\unsure{\mathrm{Cat}_{/S}}}\,\underline{\mathrm{Fib}}_W S\subset\mathrm{Cat}_{/S}
\]

\unsure{du morphisme structural|des morphismes structuraux} $X\longrightarrow S$ d'un objet, en

\diagram{\begin{tikzcd}[ampersand replacement=\&]
X\arrow[r,"i_X"]\& C(X)\arrow[r,very thick]\&S,
\end{tikzcd}}

i.e. on se donne le foncteur $C$, \uline{plus}
% archive sheet number: 242
\scanpage{243}
\mspage{193}

le morphisme fonctoriel

\[
\unsure{u|i}_{\deleted{\unsure{X|f}}}:\mathrm{id}_{\mathrm{Cat}_{/S}}\longrightarrow C
\]

\[
\text{avec}\qquad (i_X:X\mathbin{\unsure{-|\to}}C(X))\in W\qquad \forall X\ \text{dans }\mathrm{Cat}_{/S}\,;
\]

soit \deleted{\unsure{demander un}} \inserted{être encore plus exigeant et} faire varier en même temps la source et le but de $f:X\longrightarrow Y\ (=S)$, et demander un foncteur \deleted{\unsure{section}}

\[
\begin{aligned}
C:\underline{\mathrm{Fl}}(\mathrm{Cat})&\longrightarrow\mathrm{Cat},\\
f&\longmapsto C(f),
\end{aligned}
\]

avec des morphismes fonctoriels en $f$

\diagram{\begin{tikzcd}[ampersand replacement=\&,row sep=.5em]
s(f)\arrow[r,thick,"i_f"]\arrow[d,equal]\&C(f)\arrow[r,"p_f"]\arrow[d,equal]\&b(f)\arrow[d,equal]\\
X\&\overline X\&Y
\end{tikzcd}}

(\inserted{$i_{\deleted{\unsure{f}}},p_{\deleted{\unsure{f}}}$ des} flèches dans $\underline{\mathrm{Hom}}(\underline{\mathrm{Fl}}(\mathrm{Cat}),\mathrm{Cat})$),

\[
\text{avec}\qquad i_f\in W,\quad p_f\in W\text{-}\mathrm{Fib}
\]

pour toute flèche $f$ dans $\mathrm{Cat}$.
% archive sheet number: 243
\scanpage{244}
\mspage{194}

Dans les deux points de vue, dans quel sens a-t-on \uline{unicité} d'une telle construction ?

Prenons d'abord le premier point de vue, factorisation des flèches à but fixé $S$, donc question \unsure{dans} $\mathrm{Cat}_{/S}$. J'ai oublié une propriété importante, mais elle résulte de ce qui a été dit : le foncteur

\diagram{\begin{tikzcd}[ampersand replacement=\&,column sep=.9em,row sep=.55em]
X\arrow[r,mapsto]\&C(X)\,;\&\mathrm{Cat}_{/S}\arrow[rr]\arrow[dr,dashed]\&\&\mathrm{Cat}_{/S}\\
\&\&\&\underline{\mathrm{Fib}}_W S\arrow[ur,dashed,hook]\&
\end{tikzcd}}

\uline{transforme flèches} $\in W$ \uline{en \unsure{itou}}, comme on voit sur le diagramme comm.

\diagram{\begin{tikzcd}[ampersand replacement=\&]
X\arrow[r,"i_X"]\arrow[d,"f"']\&C(X)\arrow[d,"C(f)"]\\
Y\arrow[r,no head,"i_Y"]\&C(Y),
\end{tikzcd}}

où $i_X,i_Y$ sont dans $W$, donc

\[
f\in W\Longleftrightarrow C(\unsure{g|f})\in W.
\]

Soit

\[
(C',i')
\]

un deuxième foncteur, $\mathrm{Cat}_{/S}\mathbin{\unsure{-|\to}}\mathrm{Cat}_{/S}$,
% archive sheet number: 244
\scanpage{245}
\mspage{195}

muni d'un morphisme

\[
i':\mathrm{id}_{\mathrm{Cat}_{/S}}\longrightarrow C'
\]

qui est dans $W$ argument par argument, et tel que $C'$ se factorise par $\underline{\mathrm{Fib}}_W S$.

On peut comparer $C$ et $C'$ moyennant $C\circ C'$, car

\diagram{\begin{tikzcd}[ampersand replacement=\&]
\&C\circ C'\&\\
C\unsure{(}{}=C\circ\mathrm{id}\arrow[ur,very thick,"{C*i'}"]\arrow[ur,phantom,"{\deleted{\unsure{\swarrow}}}"{pos=.18,swap}]\&\&C'=\mathrm{id}\circ C'\arrow[ul,"{i*C'}"']
\end{tikzcd}}

i.e. on a, pour tt $X$ dans $\mathrm{Cat}_{/S}$

\margin{\unsure{''}}

\diagram{\begin{tikzcd}[ampersand replacement=\&]
\&C(C'(X))\&\\
C(X)\arrow[ur,"{C(i'_X)}"]\&\&C'(X)\arrow[ul,"{i_{C'(X)}}"'],
\end{tikzcd}}

ces deux flèches étant dans $W$ (celle de gauche parce que $C(W)\subset W$). Ce sont des flèches dans $\underline{\mathrm{Fib}}_W S$, et si on \deleted{\unsure{admet}} suppose que $W$ satisfait $\unsure{\mathrm{W}(4|\mathrm{WL}4|\mathrm{W}4}\,\mathrm b)$, elles induisent des $W$-équivalences en chaque fibre. Cela montre donc déjà
% archive sheet number: 245
\scanpage{246}
\mspage{196}

qu'à isos \inserted{(peut-être} non canoniques) près, le foncteur composé défini par $C$

\[
X\longmapsto C(X)\longmapsto (s\longmapsto\unsure{\mathrm{hot}|\mathrm{Hot}}_W(C(X)_s))
\]

\diagram{\begin{tikzcd}[ampersand replacement=\&,column sep=1.7em]
\mathrm{Cat}_{/S}\arrow[r]\arrow[rr,bend right=18,"{\unsure{\varphi|\phi}_C}"']\&\underline{\mathrm{Fib}}_W S\arrow[r]\&\underline{\mathrm{Hom}}(S^{\unsure{\circ|o}}\unsure{,|;}\mathrm{Hot}_W)
\end{tikzcd}}

ne dépend pas du choix particulier de $C$, i.e. deux tels foncteurs ``fibres homotopiques'' (correspondant à deux choix $C,C'$) \unsure{sont} toujours isomorphes.

Mais au lieu d'utiliser $C\circ C'$ pour définir un tel isomorphisme, on aurait pu utiliser $C'\circ C$, ou quelque autre procédé. On aimerait comprendre s'il y a unicité véritable, une caractérisation intrinsèque de ces isomorphismes $\varphi_C\xrightarrow{\sim}\varphi_{C'}$, indépendamment de telles constructions ou de telles autres.

Je présume que $C\circ C'$ est tout aussi muni de $\mathrm{id}_{\mathrm{Cat}_{/S}}\longrightarrow C\circ C'$,
% archive sheet number: 246
\scanpage{247}
\mspage{197}

et $=$\unsure{.} à priori des deux façons possibles

\diagram{\begin{tikzcd}[ampersand replacement=\&,row sep=1.4em]
\&C(C'(X))\&\\
C(X)\arrow[ur,"{C(i'_X)}"]\&\&C'(X)\arrow[ul,"{i_{C'(X)}}"']\\
\&X\arrow[ul,"i_X"]\arrow[ur,"i'_X"]\arrow[uu,dashed]\&,
\end{tikzcd}}

soit par $C(i'_X)\circ i_X$, soit par $i_{C(X)}\circ i'_X$, heureusement elles sont égales, car pour toute flèche $X\xrightarrow{i'}X'$ dans $\mathrm{Cat}_{/S}$, le carré

\diagram{\begin{tikzcd}[ampersand replacement=\&,row sep=.8em]
\&C(X')\&\\
C(X)\arrow[ur,"C(i')"]\arrow[dr,"i_X"']\&\&X'\arrow[ul,"i_{X'}"']\\
\&X\arrow[ur,"i'"]\&
\end{tikzcd}}

est commutatif (par fonctorialité de $i_X:X\mathbin{\unsure{-|\to}}C(X)$), et on applique ceci au cas de $X'=C'(X)$, $i'=i'_X$. \deleted{\unsure{Ce}}

Cette remarque nous encourage à faire un saut dans l'abstraction, et d'introduire la catégorie
% archive sheet number: 247
\scanpage{248}
\mspage{198}

\[
\begin{aligned}
\unsure{\mathcal C|\mathscr C}=\mathcal C_S=\{\ &\text{catégorie des couples }(C,i),\ \text{avec}\\
&C:\mathrm{Cat}_{/S}\longrightarrow\mathrm{Cat}_{/S}\ \text{se factorisant}\\
&\hspace{8em}\text{par }\underline{\mathrm{Fib}}_W S,\\
&i:\mathrm{id}_{\mathrm{Cat}_{/S}}\longrightarrow C\quad\text{t.q. }i_X\in W\quad\forall X\ \text{dans }\mathrm{Cat}_{/S}\ \}
\end{aligned}
\tag*{\makebox[0pt][l]{\hspace*{-\displaywidth}(\unsure{A|*})}}
\]

une flèche

\[
(C,i)\mathbin{\unsure{-|\to}}(C',i')
\]

étant une flèche

\[
C\xrightarrow{u}C'\qquad\text{(dans }\underline{\mathrm{Hom}}(\mathrm{Cat}_{/S},\mathrm{Cat}_{/S})
\]

rendant commutatif le carré

\diagram{\begin{tikzcd}[ampersand replacement=\&]
\&\mathrm{id}_{\mathrm{Cat}_{/S}}\arrow[dl,"i"']\arrow[dr,"i'"']\&\\
C\arrow[rr,"u"]\&\&C'.
\end{tikzcd}}

En d'autres termes, $\mathcal C$ est défini comme une sous-catégorie pleine de

\[
{}_{\mathrm{id}_{\mathrm{Cat}_{/S}}}\!\backslash^{\underline{\mathrm{Hom}}(\mathrm{Cat}_{/S},\mathrm{Cat}_{/S})},
\]

à savoir celle dont les objets $(C,i)$ \unsure{satisfaisant|satisfassent|satisfont} les deux conditions \unsure{(A)|de|des} plus haut (impliquant l'une et l'autre le localiseur $W$).
% archive sheet number: 248
\scanpage{249}
\mspage{199}

Ces conditions impliquent ceci

\[
\text{a)}\quad\forall X\ \text{dans }\mathrm{Cat}_{/S},\quad (u_X:C(X)\mathbin{\unsure{-|\to}}C'(X))\in W
\]

\deleted{\unsure{b)}} puisque

\diagram{\begin{tikzcd}[ampersand replacement=\&]
\&X\arrow[dl,"{i_X\in W}"']\arrow[dr,no head,"{i'_X\in W}"]\&\\
C(X)\arrow[rr,"u_X"]\&\&C'(X),
\end{tikzcd}}

et

b) Si $W$ satisfait $W(4\,\mathrm b)$, de sorte que $u_X\in W_S^f$, alors $u$ définit un \uline{iso.} entre $\unsure{\varphi|\phi}_C(X)$ et $\varphi_{C'}(X)$

\diagram{\begin{tikzcd}[ampersand replacement=\&]
s\arrow[r,mapsto]\&\unsure{\mathrm{hot}|\mathrm{Hot}}_W(C(X)_s)\arrow[d,"{\unsure{2|\wr}\ u_{X,s}}"]\\
s\arrow[r]\&\mathrm{hot}_W(C'(X)_s)
\end{tikzcd}}

\inserted{Mieux encore,} \deleted{\unsure{D}}dans ce cas, le foncteur composé

\diagram{\begin{tikzcd}[ampersand replacement=\&,column sep=.9em]
{\unsure{\mathcal C|\mathscr C}\deleted{\times\mathrm{Cat}_{/S}}}\arrow[r,very thick]\&\underline{\mathrm{Fib}}_W S\arrow[r,very thick]\&{\deleted{\unsure{\mathrm{Hot}}}\,\mathrm{HOT}^{\unsure{\mathrm{lc}|\mathrm{le}}}_W(S)}
\end{tikzcd}}

\deleted{transforme toute} donne naissance à

\[
\mathcal C\longrightarrow\underline{\mathrm{Hom}}(\mathrm{Cat}_{/S},\mathrm{HOT}^{\mathrm{lc}}_W(S))
\]

qui transforme toute flèche de $\mathcal C$ en \uline{iso}, donc se factorise par le groupoïde fondamental de $\mathcal C$

\[
\Pi\mathcal C\longrightarrow\underline{\mathrm{Hom}}(\mathrm{Cat}_{/S},\mathrm{HOT}^{\mathrm{lc}}_W(S))
\]
% archive sheet number: 249
\scanpage{250}
\mspage{200}

Quant au foncteur envisagé tantôt

\[
\unsure{\mathcal C|\mathscr C}\longrightarrow
\underbrace{\underline{\mathrm{Hom}}(\mathrm{Cat}_{/S},\underline{\mathrm{Hom}}(S^{\unsure{\circ|o}},\underline{\mathrm{Hot}}_W\deleted{\unsure{)}}))}_{\simeq\mathrm{Hom}(\mathrm{Cat}_{/S}\times S^\circ,\mathrm{Hot}_W)}
\]

\[
C\longmapsto ((X,s)\longmapsto\unsure{\mathrm{hot}|\mathrm{Hot}}_W(C(X)_s)),
\]

il se déduit du précédent en composant avec

\diagram{\begin{tikzcd}[ampersand replacement=\&,row sep=.6em,column sep=1em]
\mathrm{HOT}_W(S)\arrow[rr]\arrow[dr,dashed]\&\&\underline{\mathrm{Hom}}(S^\circ,\mathrm{Hot}_W)\\
\&\underline{\mathrm{Hom}}(\Pi S^\circ,\mathrm{Hot}_W)\arrow[ur,dashed]\&
\end{tikzcd}}

plus exactement avec le foncteur qui s'en déduit

\[
\underline{\mathrm{Hom}}(\mathrm{Cat}_{/S},\mathrm{HOT}_W(S))\longrightarrow\underline{\mathrm{Hom}}(\mathrm{Cat}_{/S},\underline{\mathrm{Hom}}(S^\circ,\mathrm{Hot}_W)\unsure{)|,}
\]

On voit poindre ici un principe plausible de canonicité \inserted{d'isomorphismes}, en ce qui concerne les foncteurs

\[
C_W:\mathrm{Cat}_{/S}\xrightarrow{\unsure{\leftarrow|\sim}}\mathrm{HOT}^{\unsure{\mathrm{lc}|\mathrm{le}}}_W(S)
\]

définis par les $C\in\mathrm{Ob}\,\mathcal C$, par la relation

\[
\begin{aligned}
\Pi\mathcal C&\approx\text{catégorie ponctuelle}\\
&\text{i.e.}\qquad\mathcal C\ \text{est }1\text{-connexe}.
\end{aligned}
\]

De plus, les arguments utilisés jusqu'ici sont de nature à tel point formelle, que les aspects spécifiques du Modéliseur $\mathrm{Cat}$
% archive sheet number: 250
\scanpage{251}
\mspage{201}

semblent n'y jouer pratiquement aucun rôle, ou un rôle très effacé.

Le carré de la p.~197 montre surtout que la catégorie $\unsure{\mathcal C|\mathscr C}$ est filtrante -- $(C,i),(C',i')$ s'envoient l'un et l'autre dans $(C'',i'')$, avec $C''=C\circ C'$, \deleted{\unsure{Alors}} $i''=(C*i')\circ i=(i*C')\circ i'$. Il y a $=$

une loi de composition dans $\mathcal C$, nullement commutative ni unitaire c'est sûr, mais sûrement associative. et \unsure{(}\deleted{\unsure{soit|si}} cette composition est notée $*$) avec des morphismes fonctoriels

\diagram{\begin{tikzcd}[ampersand replacement=\&]
{\unsure{\xi|\zeta}}\arrow[r,"{\alpha_{\xi,\xi'}}"]\&\xi*\xi'\&\xi'\arrow[l,"{\beta_{\xi,\xi'}}"']
\end{tikzcd}}

qui incarnent le caractère filtrant croissant de $\mathcal C$. Ces flèches semblent \unsure{remplacer} en quelque sorte l'unité bilatère absente, et on s'attend à ce qu'elles aient les mêmes excellentes propriétés que si on avait \unsure{une} telle unité, concernant les composés doubles $((\xi\unsure{\circ|*}\xi')\unsure{\times|*}\xi''\simeq\xi\circ(\xi'\unsure{\circ|*}\xi''\unsure{)}$ \unsure{:}
% archive sheet number: 251
